\documentclass[12pt,a4paper]{article}
\usepackage[utf8]{inputenc}
\usepackage[T1]{fontenc}
\usepackage{geometry}
\geometry{margin=2cm}

\begin{document}
\begin{center}
    \Large \textbf{PYQ for Software Enginnering}
    

    Compiled from CAT \& End-Sem Questions
\end{center}

\section{1 Marks Questions}
\begin{enumerate}
    \item Define generic software with example. \hfill [End 2026]
    \item What do you understand by the term life cycle model of the software development? \hfill [End 2026]
    \item Explain the need for SRS. \hfill [End 2026]
    \item What is the difference between module coupling and module cohesion? \hfill [End 2026]
    \item Distinguish between verification and validationdation. \hfill [End 2026]
    \item List two advatanges disadvantages of LOC. \hfill [End 2026]
    \item What are the objectives of testing? \hfill [End 2026]
    \item Discuss some of the problems that are faced during maintenance of software. \hfill [End 2026]
\end{enumerate}

\section{5 Marks Questions}
\begin{enumerate}
    \item Consider the problem of railway reservation System and design the following:
    \begin{enumerate}
        \item Problem Statement
        \item Use Case Diagram
        \item Use Cases
    \end{enumerate}
    \hfill [CAT 2026]
    
\item Consider an example of grading a student in a university. The grading is done as given below:


\begin{tabular}{|l|r|}
  \hline
  Average Marks   & Grade \\
  \hline
  90 - 100  & Exemplary performance \\
  \hline
  75 - 89 & Distinction \\
  \hline
  60 - 74 & First Division \\
  \hline
  50 - 59 & Second Division \\
  \hline
  0 - 49 & Fail \\
  \hline
\end{tabular}

The marks of any three subjects are considered for the calculation of average marks. Generate boundary value analysis test cases. Also create equivalence classes and generate test cases.
\hfill [CAT 2026]
\end{enumerate}


\section{6 Marks Questions}
\begin{enumerate}
    \item What is Software Requirement Specifiacation (SRS)? Explain the characteristics of a good Software Requirement Specification. \hfill [End 2026]
    \item \begin{enumerate}
        \item You are asked to develop an accountinf system for CUSB Gaya that replaces the existing system. Suggest the most appropiate software process development model that might be used as a bias for managing system development.Give proper justification.
        \item Under what circumstances is it benificial to construct a prototype? Does construction of a prototype always increase the overall cost of software development?
        

        \hfill [End 2026]
    \end{enumerate}
    \item Explain the types of cohesion and coupling in the context of software desing. \hfill [End 2026]
    \item Explain why the best programmers do not always make the best software managers. You may find it helpfull to base your answer on the software project managerment activities.

    \hfill [End 2026]

    \item why is accurate estimation of the effort required for completing a project difficult? Briefly explain the different effort estimations methods that are avialable. \hfill [End 2026]

    \item \begin{enumerate}
        \item Annual Change Traffic (ACT) for a software system is 13\% per year.The development effort is 600 PMs.Compute estimate for Annual Maintenance Effort (AME).If lifeline of the project is 10 years, what is the total effort of the project (using Bohem Model)?
        \item Suppose that a project was estimated to be 500 KLOC.Calculate the effort and development time for each of the three modes of Basic COCOMO model i.e.,organic, semideatched and embedded.


        \hfill [End 2026]
    \end{enumerate}
    \item 

\end{enumerate}

\section{15 Marks Question}
\begin{enumerate}
    \item What is black box testing? Explain the different types of black box testing with example.

    \hfill [End 2026]
    \item \begin{enumerate}
        \item Consider a software project with 5 task T1 - T5. Duration of the 5 task (in days) are 15,10,12,25 and 10 respectively. T2 and T4 can start when T1 is complete. T3 can start when T2 is complete. T5 can start when both T3 and T4 are complete. The project is to completed in the shortest time
        \begin{enumerate}
            \item Calculate the early time and late time for each event.
            \item Find the critical paths and the critical activities.
        \end{enumerate}

        \item What is the purpose of the data flow diagrams? Explain by constructing a context flow diagram, level -1 DFD and level -2 DFD for a library management system?

        \hfill [End 2026]
    \end{enumerate}

    \item Write short notes on any two of the following:
        \begin{enumerate}
            \item Different types of maintenance
            \item Path Testing
            \item Use Case Diagram
        \end{enumerate}
        \hfill [End 2026]
\end{enumerate}
\end{document}
